% Legends of the ERL main figures (D-100, D-102; figures/erl_fig_*.py, outputs in figures/out/).
% Shortened to at most 120 words each at the PI's request (D-107 addendum); the display details
% that the earlier legends carried (composites, contour smoothing, display grids, zoom and arrow
% rules, the Great Salt Lake level choice, symbols, projection, graticules and scale bars) are in
% the Supplementary Methods (si_erl_methods_display.tex). Order follows D-100; the audit figure
% (sections/erl_fig_audit.tex) is inserted after the shoreline map by the hook in main_erl.tex.
% Every number is a \metric{...} macro; legends spell out every acronym (ERL author guidelines).

\renewcommand{\figurename}{Figure}

\begin{figure}[!htbp]
  \centering
  \includegraphics[width=\linewidth]{../figures/out/erl_fig_then_now.pdf}
  \caption{\textbf{Three decades of retreat of the North Caspian Sea.}
  (a) Landsat true-colour composite of June--September \metric{erl.landsat.tm_1996.year}, near the high stand ($n = \metric{erl.landsat.tm_1996.n_scenes}$ scenes)\citep{usgs_l45tm_c2l2}.
  (b) The same for \metric{erl.landsat.oli_2025.year} ($n = \metric{erl.landsat.oli_2025.n_scenes}$ scenes)\citep{usgs_l89oli_c2l2}. Coloured lines, the observed shoreline at four basin levels, from the Landsat waterlines binned by level (key: level, and years of the acquisitions); dashed white line, the study domain.
  (c) Daily basin level from satellite altimetry (m, on the reference of the Database for Hydrological Time Series of Inland Waters)\citep{schwatke2015,hydrowebwl,cretaux2011}; grey bars, the two composite summers; coloured bars, the acquisitions of each level in (b). Composites and shorelines: Supplementary Methods.}
  \label{fig:erl_thennow}
\end{figure}

\begin{figure}[!htbp]
  \centering
  \includegraphics[width=\linewidth]{../figures/out/erl_fig_shorelines.pdf}
  \caption{\textbf{Where the global bathymetry grids put the shoreline.}
  (a) The North Caspian at a level of \metric{erl.hero.e3_b019.level_m}~m ($n = \metric{erl.hero.e3_b019.n_acquisitions}$ Landsat acquisitions, \metric{erl.hero.e3_b019.first_year}--\metric{erl.hero.e3_b019.last_year}): observed state (fill), observed shoreline (black) and the shorelines that GEBCO\_2026\citep{gebco2026} (blue, dashed), ETOPO 2022\citep{etopo2022,macferrin2025} (vermilion, dash-dotted) and SRTM15+\citep{tozer2019} (green, dotted) predict after their area-fitted vertical shift; boxes, the zooms.
  (b--d) Zooms on the Volga delta front, the northeastern Kazakh coast and Komsomolets Bay; the observations of the Volga delta and Komsomolets Bay are not validated (Methods).
  (e--g) The same zooms at \metric{erl.hero.e3_b029.level_m}~m ($n = \metric{erl.hero.e3_b029.n_acquisitions}$ acquisitions, \metric{erl.hero.e3_b029.first_year}--\metric{erl.hero.e3_b029.last_year}). Arrows, a distance close to the sector's median displacement; their lengths are listed in km beside each grid's line style (top to bottom: GEBCO\_2026, ETOPO 2022, SRTM15+). The zoom in (e) holds no shoreline cell near the Volga sector's median, so its arrows are longer than the median. Levels on the altimetry reference (Methods); display rules: Supplementary Methods.}
  \label{fig:erl_shorelines}
\end{figure}

\begin{figure}[!htbp]
  \centering
  \includegraphics[width=\linewidth]{../figures/out/erl_fig_gsl.pdf}
  \caption{\textbf{The same test at Great Salt Lake, Utah, where no grid holds lake-bed soundings.}
  (a) Both arms at \metric{erl.gsl.level_navd88_m}~m (North American Vertical Datum of 1988, converted from the gauges' record on the National Geodetic Vertical Datum of 1929\citep{usgs_nwis,noaa_ncat}), the level chosen by a second rule after the first picked an atypical level (Supplementary Methods): observed state, observed shoreline (black), the shorelines predicted by GEBCO\_2026, ETOPO 2022 and SRTM15+ after their best-fitting vertical shift, and that of the 3DEP control (grey), mostly 2016 lidar\citep{usgs_3dep10m,quantumspatial2017_gsl}.
  (b) Per arm and surface, the median over evaluable levels of the median distance between predicted and observed shorelines ($n = \metric{generality.great_salt_lake.gebco_2026.shape.north.n_evaluable_levels}$ levels north, \metric{generality.great_salt_lake.gebco_2026.shape.south.n_evaluable_levels} south), against the observation noise floor (grey bar: the median over levels of the per-level \metric{params.err_pct_hi}th percentiles of the distance between two interleaved halves of the observations). Filled, the surface mispredicts the arm; open, it does not. Logarithmic axis.
  (c) Heights of each grid over each arm's emerged band: GEBCO\_2026, GEBCO\_2022, ETOPO1 and SRTM15+ hold a flat surface at \metric{rev.w.g.gsl.flat_height_navd88_m}~m on \metric{rev.w.g.gsl.flat_share_pct_min}--\metric{rev.w.g.gsl.flat_share_pct_max}\% of the band, and ETOPO 2022 one at \metric{rev.g.gsl.north.etopo2022.modal_height_navd88_m}~m (north) or \metric{rev.g.gsl.south.etopo2022.modal_height_navd88_m}~m (south); dotted line, the gauge level of February 2000, when SRTM flew. Display: Supplementary Methods.}
  \label{fig:erl_gsl}
\end{figure}

\begin{figure}[!htbp]
  \centering
  \includegraphics[width=\linewidth]{../figures/out/erl_fig_agreement.pdf}
  \caption{\textbf{Where the bathymetry grids put the first metre of further fall.}
  For each cell of the study domain, the number of the five grids\citep{gebco2026,gebco2022,etopo2022,tozer2019,amante2009} that place it in the seabed emerging between the 2026 level (\metric{h3.level_2026_m}~m on the altimetry reference) and \metric{h3.scenario.1_m}~m below it: (a) measured from the observed 2026 shoreline (the observed sea at the lowest audited level, completed with the near-minimum extent of 2025 and 2026; Methods); (b) on the grids shifted by their hindcast constants. Hatched, three or more of the five grids place the cell in the band; boxes, the area that any grid and that three or more grids place there. Thick black line with a white halo, the outline of the composite baseline (the observed sea of 2025 and 2026). (c) Per grid, the share of its \metric{h3.scenario.1_m}~m band that was already dry in the extent of 2025 and 2026, as distributed (filled bars, \metric{rev.fig.f5.dry_uncorrected_pct_min}--\metric{rev.fig.f5.dry_uncorrected_pct_max}\%) and shifted (open bars, \metric{rev.fig.f5.dry_shifted_pct_min}--\metric{rev.fig.f5.dry_shifted_pct_max}\%). Five copies of one grid perturbed by its own vertical error (correlation ranges of \metric{rev.w.n.sim.range_secondary_km} and \metric{rev.w.n.sim.range_primary_km}~km) agree on \metric{rev.w.n.sim.r150km.agree_unc1_ref_p05_pct_min}--\metric{rev.w.n.sim.r50km.agree_unc1_ref_p95_pct_max}\% of their band, so low agreement here is not evidence of error beyond that dispersion (Supplementary section~\SIrevision). Display: Supplementary Methods.}
  \label{fig:erl_agreement}
\end{figure}

\begin{figure}[!htbp]
  \centering
  \includegraphics[width=\linewidth]{../figures/out/erl_fig_timeline.pdf}
  \caption{\textbf{When the sea reaches the next metres, and how much seabed each grid says it exposes.}
  (a) Basin level from satellite altimetry\citep{schwatke2015,hydrowebwl} and the published pathways continued from the 2026 level\citep{samant2023,hoseini2025}: lines, the primary pathways; shading, the range pathways (inter-model ranges and percentiles); circles, the year a primary pathway first reaches \metric{h3.scenario.1_m}~m and \metric{h3.scenario.2_m}~m below the 2026 level.
  (b) Emerged area of the study domain against the decline below the 2026 level, one curve per grid: solid, measured from the observed 2026 shoreline; dashed, the grids shifted by their hindcast constants; shading, hindcast error envelopes normalised by the decline, drawn up to the largest decline in the hindcast (\metric{rev.n.env.max_pair_decline_m}~m), beyond which the curves are uncalibrated; thin grey, the grids as distributed, an illustration of datum risk only. Pathways: Supplementary Methods.}
  \label{fig:erl_timeline}
\end{figure}
